\section{Introduction}
\label{sec:intro}

An attacker who has just landed inside a honeypot issues a single
request: \emph{show me the database credentials}. A recorder writes the
request to disk and waits. A scripted honeypot answers from a file, the
same file for every attacker, forever. A generative honeypot drafts an
answer on the spot -- and whatever that draft contains, the attacker
will believe it, act on it, and carry it into a timeline that a security
operations center will later treat as evidence. If the draft invents a
credential, that credential exists nowhere except in the attacker's
hands; if it echoes a real one, the honeypot has converted a trap into a
breach. The moment a honeypot begins to answer, it has taken on three
jobs at once: stay consistent with everything it has said so far, leak
nothing real, and remain available when the model behind it is not. A
recent survey of production deception frameworks found that no widely
deployed honeypot grounds its responses in real-time threat intelligence
or differentiates its behavior across attacker
classes~\cite{lopez2024}, and while the newest research systems manage
state and validate generative
output~\cite{li2026llmsshh,laragutierrez2026genpot}, none treats the
serving decision itself as a verified, audited security architecture
with a priced mechanism ladder (\S\ref{sec:related}). The industry's
answer to the request above is still silence.

The cost of that silence is measured elsewhere. Mean time to identify a
breach runs in months~\cite{ibm2024}, while adversary breakout time has
compressed to minutes~\cite{crowdstrike2026}. A honeypot sits directly
in the attacker's path, yet converts almost none of that proximity into
understanding: it cannot tell a credential-stuffing bot from a
hands-on-keyboard operator, and it produces none of the
tactics-and-techniques intelligence a security operations center would
need to act on the encounter.

Large language models (LLMs) promise to close this gap: prompted with a
session transcript and a threat-intelligence library, a modern LLM can
impersonate a vulnerable Linux host, improvise plausible file systems,
and steer the conversation toward the attacker's next
objective~\cite{reworr2024,wang2024,paul2025,li2026llmsshh,laragutierrez2026genpot,wang2026honeygpt}.
The promise arrives with a problem those mechanisms bound only partially:
\emph{the model can be wrong in ways a script cannot}, and a validated
response is not yet a certified one -- no recorded verdict, no bounded
failure path, no audited human gate on what gets served. A honeypot's
product is a lie, and the lie must satisfy three properties at once. It
must be \emph{consistent}: the kernel version printed on turn one cannot
contradict the one an attacker reads on turn five. It must be \emph{safe
to serve}: a model that complies too well will invent credential
material with no way to distinguish what it fabricated from what it
leaked. And it must be \emph{available}: a generator is a rented
service, and ours was -- under safety-alignment pressure from one
commercial provider, 86.7\% of generated turns exceeded the inference
budget, and response fidelity depended on how much quota the account
happened to have (deployment observation, quantified by the quota probe
log released with the stack; this figure predates the benchmark-era
watcher and is reported as an operational datum, not a benchmark
result). What no prior system does, to our knowledge, is make the
\emph{serving path itself} the security boundary -- every response
either admitted by a recorded serving-time verdict, served from vetted
deterministic content, or withheld -- and then experimentally price
which intelligent mechanisms deserve to remain inside that boundary.

Why not fix this inside the model? Prompting, refusal training, and
constrained decoding all shape probabilities, but none can certify a
particular output; alignment behavior shifts with phrasing and
pressure~\cite{ganguli2022,bai2022}, precisely the failure our quota
cliff exposed; and a deceived honeypot is itself an LLM application with
attacker-controlled context, so prompt-injection
techniques~\cite{greshake2023} apply with the attacker already in the
room. The classical system-side fixes fail differently: pre-scripting
regresses to the finite automaton of low-interaction deception, and
real-time human review cannot keep pace with an attacker who reads
faster than an analyst can vet. What survives is an asymmetry.
Generating a plausible answer to an arbitrary question is open-ended --
its output space cannot be enumerated, so no amount of testing certifies
the next draw. Checking a draft against a small, closed set of claims is
the opposite: a bounded computation whose verdict is final. A system can
exploit this by splitting the roles, letting the generator propose and a
verifier dispose against a fiction the system itself fully controls.

This paper takes the position that the right question is not
\emph{can an LLM impersonate a host} but \emph{how much intelligence can
a lie safely spend, and where must the spending stop}. Deception systems
build intelligence in rungs -- templates, classifiers, retrieval, free
generation -- and each rung has a trust price that practice pays
implicitly. We make the pricing explicit. Our answer is a verified
deception loop in which
each mechanism exists to buy one of the three properties, placed on a
single axis of certifiability (\S\ref{sec:ladder}). A fiction graph
binds every host fact to one consistent set of triples, and a verifier
checks every draft against that graph before anything reaches the wire:
consistency, enforced at serving time, not asked for at prompt time.
Baited content is drawn only from canary-registered secrets, so that any
disclosure is a measured event rather than a leak: safety by
construction. A deterministic fallback serves human-vetted transcripts
whenever generation or verification fails, so the system degrades to the
last known good answer rather than an unverified improvisation:
availability under provider failure. Beneath these sit a behavioral
profiler and a hybrid retrieval engine that ground a per-attacker
deception plan in curated CTI, giving the lie its targeting, and a
human-in-the-loop gate holds back any content class marked as requiring
approval. The loop turns generation from an act of trust into a measured
pipeline with an audit trail, and the evaluation reads the whole
apparatus as an escalation study: each rung of intelligence is measured
for what it adds over the rung beneath it, including the rungs that fold
back down when a certifiable mechanism reproduces their effect.

We evaluate whether each element of this loop contributes measurably,
using
artifact-backed benchmarks that we release alongside the paper. Four
questions structure the study:
\begin{itemize}
\item \textbf{RQ1:} Does retrieval grounding measurably change what the
honeypot detects and serves? A paired, per-scenario ablation across 90
matched rows (the first 30 released scenarios per suite) finds that
grounding lifts technique detection by $+0.317$ (95\% CI $[+0.162,
+0.480]$) and actor attribution by $+0.140$ ($[+0.083, +0.197]$) under
string-matching rubrics scored identically in both arms, while leaving
persona realism untouched. Because technique extraction operates on
retrieved documents, the raw contrast -- 35 techniques detected with
grounding against zero without it -- is a structural-dependence result
the design entails by construction (\S\ref{sec:rq1}); the measured
benefit is the rubric-scored differential, and we separate the two
explicitly.
\item \textbf{RQ2:} Can an LLM-driven honeypot remain internally consistent
and safe under sustained adversarial interaction? Across 45 targeted
re-queries the fiction graph serves zero contradictions with byte-identical
determinism, and across a 218-query live campaign the verification layer
records no unverified response served and no false-positive CTI alert.
\item \textbf{RQ3:} Can an LLM-driven honeypot produce deception an
independent judge panel finds convincing? A three-evaluator consensus
study -- each evaluator an independently prompted LLM judge persona,
distinct from the generator but sharing one judge model, so agreement is
a within-model consistency measure -- scores the full session
transcripts at $3.86/5$ (exact agreement $0.750$), with all eight
scenarios above the acceptance bar of 3 (fixed before scoring) and a
known judge bias against canary-registered bait corrected by an offline
audit of the rubric (\S\ref{sec:rq3}). No human raters participate; we
do not present this as a human evaluation.
\item \textbf{RQ4:} Does each additional rung of intelligence justify its
trust price? The escalation analysis adds each mechanism to the same
loop and prices its marginal effect: routing accuracy rises from 0.8875
to 0.9550; grounding contributes the $+0.154$ composite above; and at
the top rung, kind-aware adversarial classification corrects 12/13
assessments on 13 replayed real-log sessions and raises attacker bait
exposure by 125\% across five closed-loop replay pairs, yet adds no
measurable dwell time (a power-limited null, not an established zero
effect) and is unmeasurably subsumed by a deterministic decision floor
at that benchmark's resolution -- a directional fold-down signal that
turns the ladder from a build history into a design rule whose
generality this benchmark cannot yet certify. A partial counter-signal
is reported with equal weight: on the 13 real-log sessions the
deterministic floor misjudged every severity, while the assessed rung
corrected 12 (\S\ref{sec:jeveval}).
\end{itemize}

\noindent Our contributions are:
\begin{itemize}
\item \textbf{A verified deception architecture} (\S\ref{sec:design}) that
couples RAG-grounded generation to a fiction graph and a serving-time
verification gate, with deterministic fallback, so that every served
response is either verified against the deception state or produced from
vetted static content.
\item \textbf{A measurable grounding effect} (\S\ref{sec:evaluation}):
paired-difference bootstrap intervals over identical 90-row scenario
sets show $+0.154$ overall composite improvement from retrieval
grounding ($[+0.092, +0.221]$) under operationally defined rubrics; the
structural 35-vs-0 technique count is reported as a dependence result,
not a benefit (\S\ref{sec:rq1}); the full 218-query campaign with
grounding active engages 132 detections across 67 unique techniques.
\item \textbf{A consistency, safety, and stability benchmark} for LLM
honeypots (\S\ref{sec:evaluation}): contradiction-free serving under
adversarial re-query, next-technique steering accuracy (top-3 $0.9$), an
ablation battery (A1--A5) in which each removed mechanism costs the
system the property it was added to provide, and -- rarely reported for
LLM systems -- two independent 218-query campaigns reproducing the
composite to $+0.005$ with $0.991$ per-row pass agreement, deterministic
strata exact, with 88\% of serves attributed to verified generation and
zero fallback serves.
\item \textbf{An open, reproducible artifact} on commodity hardware
(2~vCPU, ${\sim}4$~GB RAM), with every number in this paper traceable to a
persisted result file and the judge-audit/stability chain reseeding
automatically on each new campaign run.
\item \textbf{An escalation methodology for AI in deception}
(\S\ref{sec:ladder}, \S\ref{sec:jeveval}): mechanisms placed on a
certifiability axis, each rung priced by measured marginal effect, and
the descent direction exercised on live machinery -- when a certifiable
mechanism reproduces an intelligent one's benefit, the intelligent one
comes out.
\end{itemize}

We are careful about what this paper does not claim. The deployment is
tailnet-only, so it observes no ambient internet background radiation;
the evaluation judge shares a provider with, but is a different model from, the
generator, and no human raters participate in the quality study; and the
system's availability on free-tier inference is a
constraint any practitioner should weigh. Consistency guarantees are
within-session by construction; cross-session divergence and response
latency remain open fingerprinting channels (\S\ref{sec:threatmodel}).
These limits are quantified
where they arise (\S\ref{sec:discussion}).

The remainder of the paper is organized as follows.
Section~\ref{sec:background} fixes vocabulary and the threat model.
Section~\ref{sec:design} presents the verified deception loop.
Section~\ref{sec:formal} states the serving-time invariant, the verified-serve
algorithm, and the adversarial analysis.
Section~\ref{sec:implementation} gives implementation details.
Section~\ref{sec:evaluation} reports the benchmark battery, the ablations,
the consensus study, and the stability runs; Appendix~\ref{app:artifacts}
indexes the released artifacts behind every number.
Section~\ref{sec:discussion} discusses limitations and deployment
guidance, Section~\ref{sec:related} positions the work in the deception
and LLM-security literature, and Section~\ref{sec:conclusion} concludes.
